\uuid{x3sN}
\exo7id{690}
\titre{exo7 690}
\auteur{bodin}
\organisation{exo7}
\datecreate{1998-09-01}
\isIndication{false}
\isCorrection{true}
\chapitre{Continuité, limite et étude de fonctions réelles}
\sousChapitre{Etude de fonctions}
\cpge{mpsi-mp2i.calcul-algebrique-trigonometrie.c, pcsi.calcul-algebrique-trigonometrie.c, ptsi.calcul-algebrique-trigonometrie.c, mpsi-mp2i.fonctions-valeurs-reelles-complexes.a, pcsi.techniques-calcul-fonctions.a, ptsi.techniques-calcul-fonctions.a, mpsi-mp2i.nombres-reels-suites.b, pcsi.nombres-reels-suites.a, ptsi.nombres-reels-suites.a, bcpst1.nombres-reels, bcpst1.fonctions-reelles-usuelles.a, bcpst1.limites-continuite.c}
\module{Analyse}
\niveau{L1}
\difficulte{}

\contenu{
\question{
Soit $f : \Rr \to \Rr$ définie par $f(x) = \frac{\cos x}{1+x^2}.$
Montrer que $f$ est major\'ee sur $\Rr$, minor\'ee sur $\Rr$.
D\'eterminer $\sup_{x\in \Rr} f(x)$.
}
\reponse{
Pour tout $x\in\R$ on a :
          $$ 0\leq |f(x)| = \frac{|\cos x|}{1+x^2} \leq \frac{1}{1+x^2}\leq 1. $$
Par cons\'equent, pour tout $x\in\R$, $f(x)\in [-1,1] $ donc $f$
est minor\'ee ($-1$ est un minorant), major\'ee ($1$ est un
majorant) et $\sup_{x\in\R}f(x)\leq 1$. Comme $f(0)=1$ on a
n\'ecessairement $\sup_{x\in\R}f(x) \geq 1$. Conclusion :
        $$   \sup_{x\in\R}f(x)=1.$$
}
}
